% 本文件由 LaTeX 排版 Agent 根据有效 translation.json 自动生成。
% author=Alexander of Alexandria; work_id=0622; title=论亚略异端及亚略罢免之书信

\chapter{\WorkTitle{论亚略异端及亚略罢免之书信}}

% structure: p0001 source=h1 target=omitted reason=page-title-already-rendered-by-parent

% structure: p0002 source=h2 target=section reason=relative-heading-rank; base=h2

\section{一、致\Place{君士坦丁堡}城主教\Person{亚历山大}}

致最可敬且志同道合的弟兄\Person{亚历山大}，\Person{亚历山大}在\OriginalTerm{主}{Lord}内向您致候：\par

1. 恶人的野心与贪婪之意，总是惯于对那些看似更显赫的教会暗设陷阱，借着种种借口攻击其教会的虔诚。因为他们受那在他们内运作的魔鬼煽动，贪求摆在他们前面的东西，便抛弃一切宗教顾忌，把对天主审判的畏惧践踏在脚下。关于这些事，我这受苦之人认为有必要向你的虔诚陈明，好叫你提防这类人，免得他们其中任何一个，或亲自或藉他人，胆敢踏入你的教区；因为这些术士深晓如何假借伪善以施展其欺诈，并利用满纸谎言、精心编织的书信，足能欺哄那持守纯一真诚\OriginalTerm{信德}{faith}的人。因此，\Person{亚略}与\Person{阿基莱斯}近日共谋，与\Person{科鲁图斯}的野心争胜，却变得比他更为恶劣。原来\Person{科鲁图斯}本人虽责备这些人，却为自己的恶谋找到某种口实；然而这些人，眼见他对基督的攻击，便再也不肯服从教会；反倒为自己建造贼窝，在其中不断地举行集会，日夜不停地向基督和我们倾泻诽谤。因为他们效法犹太人的样式，质疑一切虔敬的教义，便建立了一个反对基督的工厂，否认我们救主的\OriginalTerm{天主性}{Godhead}，宣讲祂只与所有其他人同等。他们收集了一切论及祂救恩计划并祂为我们而受卑抑的经文，试图由此汇集出其不虔敬的宣讲，而完全无视那些阐明祂永恒的\OriginalTerm{天主性}{Godhead}及与圣父同享无可言喻之荣耀的经文。因此，既然他们拥护犹太人和希腊人关于基督所持的不敬意见，便用尽一切办法竭力争取他们的称许；专务于那些他们向来惯于嘲笑我们的事，并且天天煽动攻击我们的骚乱和迫害。如今，他们确是将我们拖到法官的公堂，因着他们与那些被引入歧途的愚昧放荡妇女的勾结；有时，他们又将羞辱和丑名加于基督宗教，让他们的年轻女子可耻地游荡于各村各巷。更有甚者，连基督那无缝的长衣，行刑者尚且不愿分开，这些卑鄙之徒竟敢将它撕裂。\par

2. 我们确实，虽因他们诡秘行事而发现得较晚，得知他们的生活方式及不洁图谋，便经全体一致表决，已将他们逐出那朝拜基督\OriginalTerm{天主性}{Godhead}的教会之集会。但他们四处奔走攻击我们，开始投靠我们那些心意相同的同僚；表面上固然装作寻求和平与和谐，实则企图用巧言诱使其中一些人染上他们的疾病，向他们索取长篇大论的信函，好将这些信读给那些受骗者听，使他们在已陷的错误中不知悔改，并在不虔中刚硬，仿佛有主教与他们持相同意见且站在他们一边似的。再者，他们在我们中间错谬地教导和实行，以致被我们逐出的那些事，他们根本不向那些人承认，而是或默然略过，或加以遮掩，用虚假的言词和书信欺骗他们。所以，他们用看似美好且谄媚的言谈遮掩其害人的教义，蒙蔽那些心地单纯、易受欺诈之人，同时也不放过向众人诽谤我们的虔诚。于是便发生这样的事：有些人因签署了他们的信函，竟接待他们进入教会，尽管我认为最大的罪责在于那些胆敢这样做的圣职人员；因为这不仅为宗徒的规诫所不容，更是让这些人里面反对基督的魔鬼作为借此越发煽旺。因此，亲爱的弟兄们，我毫不迟延地振作起来，向你们揭示这些人的不信实；他们说曾有一段时间，\OriginalTerm{圣子}{Son of God}并不存在；又说那先前不存在的，后来才进入存在，成为如此，即当祂最终受造时，如同每个人惯于出生一样。因为，他们说，天主从非有之物中造了万有，甚至将\OriginalTerm{圣子}{Son of God}也包含在一切理性与无理性之物的创造之中。由此他们进而推论，祂具有可变的本性，既能行德亦能行恶。一旦假设祂是\Quote{从非有之物而来}，他们就推翻了\WorkTitle{圣经}关于祂永恒性的经文，那些经文意指\OriginalTerm{智慧}{Wisdom}与\OriginalTerm{圣言}{Word}的\OriginalTerm{不变性}{immutability}和\OriginalTerm{天主性}{Godhead}，而智慧与圣言就是基督。\par

3. 因此，这些恶人说，我们也能像祂一样成为天主的子女。因为经上记载说：\BibleQuote{我养育儿女，将他们养大。}\BibleRef{依 1:2}但当有人向他们提出接下来的话\BibleQuote{他们竟背叛我}，这实在不适用于救主的本性，因救主具有不变的本性，他们便抛弃一切宗教敬畏，说什么天主因为预知且预见祂的\OriginalTerm{圣子}{Son}不会背叛祂，就从众人中拣选了祂。因为他们说，祂拣选祂，并非因祂本性具有什么超越其他儿子们的特别之处，因为按他们所言，没有人天生就是天主的儿子；也非因祂有什么特有的品质；而是天主拣选了一个本性可变的人，因着祂品行的谨慎与修持，丝毫不转向邪恶；这样，倘若\Person{保禄}和\Person{伯多禄}也曾为此努力，他们的子性便与祂的无异了。为证实此狂谬教义，他们玩弄\WorkTitle{圣经}，引证\WorkTitle{圣咏}论及基督的话：\BibleQuote{你爱护正义，憎恨邪恶，为此天主，你的天主，以喜乐之油傅你，胜过你的同伴。}\par

4. 然而，圣子并非\Quote{从虚无中}造成，也没有一个\Quote{时间当祂不存在}，圣史\Person{若望}已充分说明，他论及祂时如此写道：\BibleQuote{在父怀里的独生子。}\BibleRef{若 1:18} 这位神圣的导师为要表明圣父与圣子是彼此不可分离的两个，便说祂是\Quote{在父怀里}。同样，天主的圣言并非属于\Quote{从虚无中}受造之物之列，同一\Person{若望}说：\BibleQuote{万物是藉着祂而造成的。} 他明示了祂固有的位格，说：\BibleQuote{在起初已有圣言，圣言与天主同在，圣言就是天主。万物是藉着祂而造成的；凡受造的，没有一样不是由祂而造成的。}\BibleRef{若 1:1} 因为，倘若万物是藉着祂而造成的，那么那位将存在赋予受造物的，怎么自己曾一度不存在呢？造物的圣言，不能界定为与受造物同性同体；因为祂本就在起初，万物都是藉着祂并由\Quote{从虚无中}而造成。再者，那自有者似乎与\Quote{从虚无中}造成的万物相反且远离。因为那确实表明圣父与圣子之间毫无间隔，即使思想也不能想象其间有任何距离。然而，世界是由\Quote{从虚无中}造成的，这表明质体的本原有其较晚的起源，因为宇宙由圣父通过圣子得到了这样一种本质。因此，当至虔诚的\Person{若望}默观天主圣言的本质，发现它极为超越，远非一切受生者所能理解，他便认为不宜谈论其生成与受造；不敢以论及受造物的措辞来指称造物主。并非圣言是\OriginalTerm{非受生}{unbegotten}的，因为唯有圣父是非受生的，而是因为独生子的\OriginalTerm{无可言喻的本体}{subsistence}超越了圣史们、或许也超越天使的敏锐理解。\par

5. 因此，我认为凡胆敢探究超越这些事之外的，不应算作虔诚者，因他不听从这话：\Quote{不要探求那太难的事，也不要搜寻超过你能力的事。} 因为若许多远不及此的事理还有人不能领会，如宗徒\Person{保禄}所说：\BibleQuote{天主为爱祂的人所准备的，是眼所未见，耳所未闻，人心所未想到的。}\BibleRef{格前 2:9} 天主也曾对\Person{亚巴郎}说，\Quote{他不能数尽星辰；}\BibleRef{创 15:5} 以及那节经文：\Quote{谁能数清海沙，数尽雨滴？}\BibleRef{德 1:2} 那么，除非癫狂，谁还能过于好奇地探究天主圣言的\OriginalTerm{本体}{subsistence}呢？关于这事，预言之神说：\Quote{有谁会描述他的身世呢？}\BibleRef{依 53:8} 我们的救主自己，那位祝福世界一切支柱者，也设法解除他们对此事的求知，说，理解这事完全超出他们的本性，这至神圣奥迹的知识只属于圣父。祂说：\BibleQuote{除了父外，没有人认识子；除了子和子所愿意启示的人外，也没有人认识父。}\BibleRef{玛 11:27} 关于这事，我认为圣父也说过：\Quote{我的秘密归于我，也归于属我的。}\par

6. 如今，认为圣子是从虚无中造成的，且在时间中有了存在，这是疯狂的想法，\Quote{从虚无中}这说法本身就表明了，尽管这些愚昧人不明白自己言语的疯狂。因为\Quote{不存在}这一表述必须要么在时间里、要么在某个时期的期限内来衡量。但若果真\BibleQuote{万物是藉着祂而造成的}，那么可以确定，每个时期、时间、所有空间，以及那包容\Quote{不存在}的\Quote{何时}，都是藉着祂造成的。而如果说那位造了时间、时期和季节（那\Quote{不存在}即交织于其中）者曾一度不存在，这岂不荒谬吗？因为声称万物之因的那一位在时间上后于那事物的起源，是毫无理智且极为无知的。因为照他们的说法，他们说圣子尚未由圣父造成的那个时间段，竟先于创造万物的天主智慧，而圣经若称祂为\Quote{一切受造物的首生者}，照他们就是虚假的。与此相符，那威严发言的\Person{保禄}论及祂说：\BibleQuote{立祂为万有的承继者，并藉着祂造成了宇宙。因为在天上和在地上的一切，可见的与不可见的，无论是上座者，或是宰制者，或是率领者，或是掌权者，都是在祂内受造的：万有都是藉着祂，并且是为了祂而受造的；祂在万有之先。}\BibleRef{哥 1:16}\par

7. 因此，既然这种\OriginalTerm{从虚无中创造的假说}{hypothesis of a creation from things which are not}显得极为不虔，就必须肯定圣父永远是圣父。但祂是圣父，正因为圣子永远与祂同在，因着圣子，祂才被称为圣父。所以，既然圣子永远与祂同在，圣父就是永远完美的，丝毫不缺任何善；祂不是在时间内，不是经过一段间隔，也不是从虚无中，生了自己的\OriginalTerm{独生子}{only-begotten Son}。那么，怎能说不虔的是，宣称天主的智慧曾一度不存在---这智慧论及自己时说：\Quote{我曾与祂同在，创造万物；我是祂的喜悦；}，或说天主的大能曾一度不存在，或说圣言曾在某个时刻残缺，或说其他使圣子为人所知、圣父得以彰显的事物曾有所欠缺？因为谁否认\OriginalTerm{光荣的光辉}{brightness of the glory}曾存在，也就除去了这光辉所源自的原初之光。如果\OriginalTerm{天主的肖像}{image of God}不是永远存在，那么显然，肖像所描绘的那一位也就不是永远存在。再者，若说\OriginalTerm{天主本体的印记}{character of the subsistence of God}未曾存在，那么由它完美表达的那一位也就被废除了。由此可见，我们救主的\OriginalTerm{子性}{Sonship}与其余众子的子性毫无共同之处。因为正如已经指明的，祂那不可言喻的本体，以无可比拟的卓越超越祂所赋予存在的其他一切，同样，祂那依照圣父\OriginalTerm{天主性}{Godhead}之本性的子性，也以无法言喻的卓越，超越那些被祂收为\OriginalTerm{义子名分}{adoption}之人的子性。因为祂确实具有\OriginalTerm{不可变的本性}{immutable nature}，各方面全美，毫无欠缺；而他们，由于无论如何都会改变，需要祂的扶助。天主的智慧能有什么进步？真理本身与天主圣言能有什么增长？生命与真光在什么方面能变得更好？如果是这样，那么说智慧竟会愚拙，天主的大能竟与软弱结合，理性竟被无理蒙蔽，黑暗竟与真光混合，该是何等违背本性啊！关于这一点，宗徒说：\BibleQuote{光明与黑暗有什么相通？基督与贝里雅耳有什么协调？}\BibleRef{格后 6:14}撒罗满也说，人不可能理解\Quote{蛇在岩石上的道}，按保禄的意见，这岩石就是基督。但人类和天使，都是祂的受造物，领受了祂的祝福，为能进步，在德行与法律的诫命中操练自己，好不犯罪。因此，我们的主，由于本性是圣父的儿子，受到万有的朝拜。但他们，却脱去\OriginalTerm{奴役的精神}{spirit of bondage}，借着勇敢的行为和进步，领受了\OriginalTerm{义子的精神}{spirit of adoption}，蒙那位本性为子者的祝福，成为义子。\par

8. 圣保禄宣告了祂那\OriginalTerm{本有的、独特的、本性的、卓越的子性}{proper and peculiar, natural and excellent Sonship}，他论及天主说：\BibleQuote{祂没有怜惜自己的儿子，反而为我们，}我们原不是祂的\OriginalTerm{本性的儿子}{natural sons}，\BibleQuote{把祂交出了。}\BibleRef{罗 8:32}为了将祂与那些并非真正儿子的人区分开，他说祂是自己的儿子。在福音中我们读到：\BibleQuote{这是我的爱子，我所喜悦的。}\BibleRef{玛 3:17}此外，在\WorkTitle{圣咏集}中，救主说：\BibleQuote{上主对我说：\InnerQuote{你是我的儿子。}}这里，为显示祂是真确的、纯正的子，祂指明除了祂自己，没有别的纯正的子。而这话又是什么意思：\BibleQuote{在晨星之前，由母胎中我生了你}？这不是分明指示出父生养的\OriginalTerm{本性子性}{natural sonship of paternal bringing forth}吗？祂获得这子性，不是凭借品行的小心培育，不是靠着德行的操练与增长，而是因着本性的特属。因此，圣父的\OriginalTerm{独生子}{only-begotten Son}确有\OriginalTerm{无缺陷的子性}{indefectible Sonship}；但\OriginalTerm{理性之子的义子名分}{adoption of rational sons}，并非生来属他们，而是因着他们生活的正直和\OriginalTerm{天主的自由恩赐}{free gift of God}，为他们预备的。并且这义子名分是\OriginalTerm{可变的}{mutable}，正如圣经所承认的：\BibleQuote{当天主的儿子们看见人的女儿们美丽，随意选取，作为妻子，}\BibleRef{创 6:2}等等。在另一处：\BibleQuote{我抚养、教育了子女，他们却背叛了我。}\BibleRef{依 1:2}如我们在天主借依撒意亚先知所说的话中见到的。\par

9. 亲爱的弟兄们，虽然我还能说许多话，我特意省略了这一切，因为我认为长篇大论地向与我同心合意的教师们重提这些事，是不胜负担的。你们自己蒙受了天主的教导，也不是不知道，这最近兴起反抗教会敬虔的学说，正是\Person{厄彼庸}与\Person{阿耳忒玛斯}的学说；它无非是模仿安提约基主教\Person{撒摩撒达的保禄}，他经过全体主教在各处的审判和磋商，被隔绝于教会之外。\Person{路基安}继他之后，多年与三位主教的共融隔绝。而今，最近有从他们的不敬之渣滓中吸取了糟粕的人，——他们隐藏的苗裔，\Person{亚略}和\Person{阿基利斯}，以及那群与他们的邪恶聚在一起的人——在我们中间兴起，教导这从非有之物创造的道理。在\Place{叙利亚}有三名主教，因与他们同心，在某些方式下被祝圣，从而煽动他们行更恶的事。但关于这些事的判断，留待你们审断。因为他们记住那些论及祂救赎苦难、卑微、受审，以及他们所谓祂的贫穷，简言之，就是救主为我们缘故所忍受的一切事，就把这些作为论据，来反驳祂的至高永恒天主性。但对于那些意指祂本然的荣耀、尊贵，以及与父同在的言语，他们却忘记了。例如这话：\BibleQuote{我与父原是一体}，\BibleRef{若 10:30}。这确实是主说的，并非宣布祂自己就是父，也不是表明两个位格是一体；而是说，父的子最精确地保存着父的明确肖像，因为祂在各方面都本然地铭刻着父的相似性，是父毫无差别的肖像，是原始模范的明确形像。因此，主也向当时渴望看见祂的\Person{斐理伯}，充分显示了这一点。因为当他说\BibleQuote{将父显示给我们}，\BibleRef{若 14:8}，主回答说：\BibleQuote{谁看见了我，就看见了父}，因为父正是借着那纯洁而生动之天主肖像的镜子，显现了祂自己。圣人们在\WorkTitle{圣咏集}中所说的也与此相似：\BibleQuote{我们在你的光中，必得见光}。因此，凡尊敬子的，也尊敬父；这言之有理，因为他们胆敢攻击子的一切不敬之言，都是指向父的。\par

10. 亲爱的弟兄们，在这之后，如果我提出他们攻击我和我们最虔诚平信徒的虚假诽谤，那有什么可惊奇的呢？因为那些列阵反对基督天主性的人，毫不犹豫地发出攻击我们的忘恩负义的狂言。他们既不愿意任何古人被比作他们，也不允许任何我们从幼年时就师从的教师被置于与他们同等的地位。不，他们甚至认为现今在我们同僚中没有任何人达到哪怕中等程度的智慧；他们自夸是唯一智慧并舍弃世物的人，是教义的独一发现者，只有他们才得到了以前天下无人曾想到的启示。哦，亵渎的傲慢！哦，不可测度的疯狂！哦，与疯狂者相称的虚荣！哦，那在他们不洁灵魂中扎根的\Person{撒殚}的骄傲！古代\WorkTitle{圣经}的宗教清晰性未能使他们羞愧，我们同僚间关于基督的一致教义也未能抑制他们对抗基督的胆大妄为。甚至连时刻伺机反对圣子的亵渎之言的恶魔，也忍受不了他们的不敬。\par

11. 现在让我们按己力提出这些论点，以驳斥那些对基督一无所知、彷佛在尘土中翻滚攻击基督的人；他们着手诽谤我们对祂的虔敬。因为那些捏造愚蠢寓言的人说：我们这些对那些说基督从非有之物而来之人的亵渎基督、不敬、不合圣经的谬论，避而远之，从而断言有两个\OriginalTerm{不受生者}{unbegotten}。因为他们无知地断言，必然要说二者之一：要么祂是从非有之物而来，要么就有两个不受生者；这些无知的人不晓得，\OriginalTerm{不受生的父}{unbegotten Father}，与祂从非有之物中创造的万有——无论是理性的还是非理性的——之间有多大的差别。在这两者之间，居于中间位置的，是天主唯一\OriginalTerm{受生的}{begotten}性体，即圣言，父借以从无中造成万物，祂由真父自身所生。正如主在某处亲自见证的，说：\BibleQuote{凡爱生祂者的，也必爱从祂所生者}，\BibleRef{若一 5:1}。\par

12. 关于这位，我们所信的，一如宗徒的教会所信的。我们信仰一位\OriginalTerm{非受生}{unbegotten}的圣父，祂的存在不来自任何人，祂是不可改变的、永恒不变的，祂是常存的，不容有任何增减；祂将法律、先知和福音赐给了我们；祂是圣祖、宗徒及所有圣人的主。并信仰唯一的主耶稣基督，\OriginalTerm{天主的独生子}{only-begotten Son of God}；祂不是由虚无中受生，而是由圣父；不是按照形体，如\Person{撒伯里乌}与\Person{瓦伦廷}所设想的分割或分离，而是按照某种不可言喻、无法言说的方式，正如上面引用的先知的话：\BibleQuote{谁会讲述祂的世代？}\BibleRef{依 53:8} 既然祂的\OriginalTerm{存在}{subsistence}是任何受生的本性都无法探究的，正如圣父不能被任何人探究一样；因为理性本性的受造物无法认识祂由圣父所得的神性的受生。然而，那些受真理之神感动的人，无需由我学习这些事，因为在我们耳中回荡着基督先前就此所说的话：除非子以外，没有人认识父；除非父以外，也没有人认识子是谁。\BibleRef{玛 11:27} 我们已得知，祂与圣父同样地不可改变、永恒不变，毫无欠缺，是完美的圣子，与圣父相似；祂只在这一点上低于圣父，即祂不是非受生的。因为祂是圣父的完全的真像，与祂毫无差异。因为显然，祂是那完全包含一切、藉此宣告出极大\OriginalTerm{相似性}{similitude}的真像，正如主亲自教导我们的，当祂说：\BibleQuote{父比我大。}\BibleRef{若 14:28} 为此我们相信，圣子出自圣父，永远存在。\BibleQuote{因为祂是天主光荣的辉耀，是\Emph{圣父}\OriginalTerm{位格}{person}的真像。}\BibleRef{希 1:3} 但是，谁也不要将那\Emph{永远}一词理解为祂是非受生的，就像那些心灵盲目的人所想象的那样。因为无论是\Quote{祂曾在}、\Quote{永远}，还是\Quote{在万世以前}这些词语，都不等同于非受生。但人的心智也不能用任何其他词语来表示非受生。因此，我想你们是明白的，我也信赖你们在一切事上的正确意向，因为这些词语根本不表示非受生。因为这些词语似乎仅仅表示时间的延续，但\OriginalTerm{天主性}{Godhead}，以及仿佛独生子的亘古，它们不能恰当地表示；而是圣人们使用了它们，各人按自己的能力寻求表达这奥秘，请求听者宽容，并提出合理的理由说，我们只能达到这一步。但若有人期望从可朽的唇舌中听到某种超越人能力的话语，说那些部分了解的已被废去，那么显然，\Quote{祂曾在}、\Quote{永远}、\Quote{在万代以前}这些词，远不及他们所希望的。而且无论采用什么词，都不等同于非受生。因此，对于非受生的圣父，我们确应保存祂独有的尊严，承认没有谁能成为祂存在的原因；但圣子则应享有祂应得的尊荣，如我们所说，赋予祂由圣父无始的出生，并给予祂\OriginalTerm{钦崇}{adoration}，只是要虔诚而恰当地对祂使用\Quote{祂曾在}、\Quote{永远}、\Quote{在万世以前}这些词；决不否认祂的天主性，而是归于祂一种相似性，这种相似性在各方面完全符合于圣父的\OriginalTerm{真像}{Image}与\OriginalTerm{典范}{Exemplar}。但我们必须说，唯独圣父具有非受生的特性，因为救主亲自说过：父比我大。\BibleRef{若 14:28} 除了关于圣父与圣子的虔敬看法外，我们宣认唯一圣神，如圣经所教导的；祂曾引导了旧约的圣者和那称为新约的神圣导师。此外，还宣认一个唯一、至公、从宗徒传下来的教会，她永远不能被摧毁，即使全世界都试图与她争战；但她在一切异端者反对她的最邪恶叛乱中，都将得胜。因为她的新郎曾说过话来坚固我们的心：\BibleQuote{你们放心，我已战胜了世界。}\BibleRef{若 16:33} 此后，我们知道死者的复活，其初果就是我们的主耶稣基督，祂确实地、而非仅在外表上，从圣母玛利亚，天主之母，取了一个身体；在世界的末了，祂来到人类中间，为除免罪过，被钉在十字架上，死了；然而祂的天主性并未因此受到任何损害，祂从死者中复活，被接升天，坐在尊威的右边。\par

13. 这些事在这封书信中只写了一部分，因为我认为逐一详尽写出过于繁重，正如我之前说过的，因为这些事逃不过你们虔诚的留意。我们这样教导，这样宣讲。这些就是教会的宗徒教义\OriginalTerm{宗徒教义}{apostolic doctrines}，我们为此也不惜一死，认为那些强迫我们背弃这些教义的人微不足道，即使他们用酷刑强迫我们，我们也不舍弃在这些教义上的希望。亚略和阿基勒斯反对这些教义，以及那些与他们一起成为真理的敌人者，已被逐出教会，因为他们与我们的神圣教义相悖，正如圣保禄所说的：\BibleQuote{但即使是我们，或是从天上来的天使，若给你们宣讲的福音，与我们所宣讲的不同，应受诅咒。} \BibleRef{迦 1:8} 又说：\BibleQuote{若有人传讲异端道理，不遵从我们的主耶稣基督的健全道理，以及那合乎虔敬的教训，他是傲慢的，一无所知，} \BibleRef{弟前 6:3} 等等。所以，这些已被弟兄团体绝罚的人，你们不可接待任何人，也不可容受他们所说的或所写的东西。因为这些诱惑者总是在说谎，他们从不说真话。他们周游各城，无非是想藉着友谊的伪装与和平的名义，用他们的伪善和谄媚，传递书信，藉此欺骗一些 \BibleQuote{轻浮的妇女，满身罪债，被他们掳去，} \BibleRef{弟后 3:4} 等等。\par

14. 所以，这些胆敢这样反对基督的人；他们一部分公然嘲弄基督宗教，一部分试图在审判台前诋毁告发其信奉者；他们在和平时，尽其所能，煽起对我们的教难；他们削弱了基督那不可言喻的受生奥迹\OriginalTerm{基督的受生奥迹}{ineffable mystery of Christ's generation}；我说，亲爱的志同道合的弟兄们，你们要厌恶地远离这些人，同我们一起反对他们疯狂的胆大妄为；正如我们的同工们所做的那样，他们义愤填膺，既给我们写了信反对这些人，又签署了我们的信。这信我也托付我的儿子执事阿皮翁带给你们，信中的签名有些来自全埃及和底比斯，有些来自利比亚和五城区。还有一些来自叙利亚、吕西亚、旁非利亚、亚细亚、卡帕多细亚，以及其他邻近省份。仿效这些，我相信我也将收到你们的信。虽然我已预备了许多帮助来治愈那些受伤的人，但这是特别为治疗受他们欺骗的群众所设计的药方，使他们遵从我们同工们的共同意见，从而赶紧回归悔改。你们要彼此问安，并问你们那里的弟兄安。我祈祷你们在主内强壮，亲爱的，也希望我能因你们对基督的爱而受益。\par

% structure: p0018 source=h2 target=section reason=relative-heading-rank; base=h2

\section{二、公函}

致各地天主教会内我们至爱和至可敬的同工，亚历山大在主内问候你们：\par

1. 既然天主教会的身体是唯一的，且圣经命令我们应当持守同心合意与和平的联系，因此理应我们互相通告各自所行的事；好使若一个肢体受苦或欢乐，我们大家一同受苦或欢乐。那么，在我们的教区内，不久以前，出现了不法之徒和反对基督的人，他们教导人背教；这事实在理所当然地可以怀疑并称之为反基督的先驱。我曾想将这事情秘而不宣，或许这样恶事只会在异端首领们中间自己消灭，而不致扩散到其他地方并污染那些较为单纯者的耳朵。但是，因为尼科美底亚现任主教优西比乌，自以为所有教会事务都取决于他，因为他离开了贝里图斯而觊觎尼科美底亚人的教会，且未受任何惩罚，就得以管辖这些背教者，并四处写信赞扬他们，想方设法要引诱某些无知的人偏向这最可耻的、反对基督的异端；所以我有必要不再沉默，因为知道律法上的话，而要向你们大家宣布，好使你们认识那些已经背教的人，也认识他们异端谬理的可悲言论；并且如果优西比乌写信来，不要理会他。\par

2. 因为他想藉着他们的帮助，重振他那久已隐藏，但曾一度保持沉默的内心旧恶，他假装为他们写信，但却用自己的行动证明，他是在为自己的缘故而尽力这样做。现在，背弃教会的人乃是这些：亚略、阿基勒斯、艾塔勒斯、卡波尼斯、另一个亚略、萨尔马特斯，曾是司铎；优佐依乌斯、路基乌斯、尤利乌斯、梅纳斯、赫拉迪乌斯、盖恩斯，曾是执事；和他们一起的还有塞孔都斯和特奥纳斯，这两人曾被称为主教。他们所捏造并述说，与圣经之意相悖的言论如下：——\par

\Quote{天主并不是常为圣父；曾有一段时期天主还不是圣父。天主的圣言并非永恒，而是从\InnerQuote{从非有之物}受造的；因为那为天主的一位，从非有之物中造成了那不存在的；因此曾有一段时期他并不存在。因为圣子是受造之物，是被造的：他在\OriginalTerm{性体}{substance}上也不相似圣父；他既不是圣父真实而本性的圣言，也不是圣父真实的智慧；他反而是受造之物中的一员。他被误称为圣言和智慧，因为他自身是由天主内真正的圣言、和那在天主内的智慧所造成的，天主怎样藉此智慧创造了其他万物，也怎样创造了他。因此，他的本性就是可变易的，与其他有理性的受造物一样。圣言也与天主的性体格格不入、彼此分离。圣父为子也是不可言喻的；因为圣言既不完全确切地认识圣父，也不能完全地看见圣父。因为圣子确实也不知道他自身的性体究竟为何。他是为了我们而被造的，好让天主藉着他有如藉着工具来创造我们；若非天主愿意造我们，他也就不会存在。有人问他们说，天主子能否像魔鬼一样变化；他们竟不惮于回答说：能；因为他是受造之物，所以具有可变的本性。}\par

3. 既然亚略一伙的人讲这些话，并且恬不知耻地坚持，我们同埃及及利比亚的主教们，数目近百人，聚集一处，已将这些人及其附从者一并施以绝罚。但欧瑟比一伙的人却接纳了他们，竭力要把虚妄与真理，亵渎与虔敬混杂起来。但他们不会得逞；因为真理终归得胜，\BibleQuote{光明与黑暗不能相通，基督与贝里雅耳不能相合。}\BibleRef{格后 6:14} 谁曾听过这样的事呢？现在听了的人，谁不惊讶，谁不捂住耳朵，以免这些污秽的话沾染自己呢？谁听了若望说，\BibleQuote{在起初已有圣言，}\BibleRef{若 1:1} 能不谴责那些说曾有一段时间他并不存在的人呢？谁听了福音中的这些话：\BibleQuote{独生子；}\BibleRef{若 1:18} 以及 \BibleQuote{万物是藉着他而造成的，}\BibleRef{若 1:3} 能不憎恨那些宣称他是受造物之一的人呢？因为，他怎可能是藉他而造成的万物之一呢？或者，若他果真是受造之物，他怎能如他们所说，被列于其他万物中，却又是独生子呢？当圣父说：\Quote{我的心灵涌溢美善的圣言；} 又说：\Quote{在晨星之前，我由胎中生了你；} 他怎能是从非有之物造成的呢？再说，他是圣父完美的肖像和光辉，他怎能与圣父的性体不相似呢？他曾说：\BibleQuote{谁看见了我，就是看见了父。}\BibleRef{若 14:9} 如果圣子是圣言、是天主的智慧与理智，又怎能曾有一段时间他并不存在呢？这无异于说，曾有一段时间天主没有理智与智慧。他亲口说：\BibleQuote{我在父内，父也在我内，}\BibleRef{若 14:10} 又说：\BibleQuote{我与父原是一体；}\BibleRef{若 10:30} 又藉先知说：\BibleQuote{我是上主，我决不改变。}\BibleRef{拉 3:6} 他怎能是可变的、易变的呢？虽然有一句话是指圣父自己，但如今用在圣言身上更为贴切，因为当他降生成人时，他并未改变；正如宗徒所说：\BibleQuote{耶稣基督昨天、今天、直到永远，常是一样。}\BibleRef{希 13:8} 是谁诱导他们说，他是为了我们的缘故而被造的；虽然保禄说：\BibleQuote{万物都出于他，并藉他而存在？}\BibleRef{希 11:10}\par

4. 至于他们说圣子并不完全认识圣父的亵渎之言，我们不必惊奇：因为他们既决意与基督为敌，就也攻击他的这句话：\BibleQuote{正如父认识我，我也认识父一样。}\BibleRef{若 10:15} 因此，如果圣父只是部分地认识子，那么显然子也不完全认识圣父。但如果这样说是不虔敬的，而且圣父完全认识子，那么很明显，就如圣父认识他自己的圣言，圣言也认识他自己的圣父，他是圣父的圣言。\par

5. 我们讲述这些事，并阐明天主圣经，经常驳斥他们。但他们像变色龙一样，改变自己的主张，企图将那句话用在自己身上：\BibleQuote{恶人一来，轻视也到。}\BibleRef{箴 18:3} 在他们之前，确有许多异端存在，它们因胆大妄为，已陷入疯狂。但这些人用尽一切言论，企图废除基督的天主性，他们因为更接近敌基督，倒使先前那些异端显得像个义人了。因此，他们被教会开除教籍、施以绝罚。虽然我们为这些人的沉沦感到悲伤，尤其因为他们曾一度学习了教会的教义，如今却倒退回去；但我们并不惊奇；因为同样的事也发生在依默乃和非肋乃身上，\BibleRef{弟后 2:17} 在他们以前，还有犹达斯，他虽然跟随救主，后来却成了叛徒和背教者。此外，关于这些人，我们不是没有预先的警告，因为主曾预言说：\BibleQuote{你们要谨慎，不要受欺骗！因为将有许多人，假冒我的名字而来，说：我就是基督；又：时期近了。你们切不可跟随他们。}\BibleRef{路 21:8} 保禄也从救主领受了这些事，写道：\BibleQuote{在末后的时期，有些人要离弃信德，听从诱人的邪神和魔鬼的训言。}\BibleRef{弟前 4:1}\par

6. 因此，我们的主和救主耶稣基督既亲自这样劝勉我们，又藉祂的宗徒将这些事指示了我们；我们既亲耳听闻了他们的不敬，正如我已说过的，便一致将这样的人予以\OriginalTerm{绝罚}{anathematized}，并宣布他们为\OriginalTerm{天主教会}{Catholic Church}和信仰的局外人。亲爱的、最可敬的同工们，我们已将这事告知你们的虔敬，好使你们若有人胆敢擅自到你们那里，你们不要接待他们中的任何人，也不要信任\Person{优西比乌}或其他任何为他们写信的人。因为我们作为基督徒，理当转离所有说话或思想敌对基督的人，视他们为天主的仇敌和灵魂的毁灭者，并且\BibleQuote{连问安也不可，以免有分於他的恶行}，正如真福\Person{若望}所吩咐的。请问候与你们同在的弟兄们。与我同在的弟兄们也问候你们。\newline{}\Signature{亚历山大城的众长老谨启}\par

\Signature{我，\Person{高路多}，长老，赞同所写之事，并赞同罢黜\Person{亚略}及与之同犯不敬之罪的一干人等。}\par

\Signature{\Person{亚历山大}，长老}\newline{}\Signature{\Person{阿波克拉提翁}，长老}\newline{}\Signature{\Person{狄奥斯科若斯}，长老}\newline{}\Signature{\Person{阿加图斯}，长老}\newline{}\Signature{\Person{尼美修斯}，长老}\newline{}\Signature{\Person{狄奥尼修斯}，长老}\newline{}\Signature{\Person{隆古斯}，长老}\newline{}\Signature{\Person{西尔瓦努斯}，长老}\newline{}\Signature{\Person{优西比乌}，长老}\newline{}\Signature{\Person{佩路斯}，长老}\newline{}\Signature{\Person{阿皮斯}，长老}\newline{}\Signature{\Person{亚历山大}，长老}\newline{}\Signature{\Person{普罗特里乌斯}，长老}\newline{}\Signature{\Person{保路斯}，长老}\newline{}\Signature{\Person{尼拉若斯}，长老}\newline{}\Signature{\Person{基若斯}，长老}\newline{}\Signature{\Person{阿摩尼乌斯}，执事}\newline{}\Signature{\Person{安比提亚努斯}，执事}\newline{}\Signature{\Person{加约}，执事}\newline{}\Signature{\Person{马喀里乌斯}，执事}\newline{}\Signature{\Person{皮斯图斯}，执事}\newline{}\Signature{\Person{亚历山大}，执事}\newline{}\Signature{\Person{狄奥尼修斯}，执事}\newline{}\Signature{\Person{亚大纳削}，执事}\newline{}\Signature{\Person{阿加彤}，执事}\newline{}\Signature{\Person{优美尼斯}，执事}\newline{}\Signature{\Person{波利比乌斯}，执事}\newline{}\Signature{\Person{阿波罗尼乌斯}，执事}\newline{}\Signature{\Person{欧林匹乌斯}，执事}\newline{}\Signature{\Person{特奥纳斯}，执事}\newline{}\Signature{\Person{阿夫托尼乌斯}，执事}\newline{}\Signature{\Person{玛尔库斯}，执事}\newline{}\Signature{\Person{亚大纳削}，执事}\newline{}\Signature{\Person{科莫杜斯}，执事}\newline{}\Signature{\Person{马喀里乌斯}，执事}\newline{}\Signature{\Person{塞拉皮翁}，执事}\newline{}\Signature{\Person{尼鲁斯}，执事}\newline{}\Signature{\Person{保路斯}，执事}\newline{}\Signature{\Person{若马努斯}，执事}\newline{}\Signature{\Person{贝特鲁斯}，执事}\par

\Place{马雷奥蒂斯}的众长老。\par

\Signature{我，\Person{阿波罗尼乌斯}，长老，赞同所写之事，并赞同罢黜\Person{亚略}及与之同犯不敬之罪的一干人等。}\par

\Signature{\Person{英革尼乌斯}，长老}\newline{}\Signature{\Person{狄奥斯科若斯}，长老}\newline{}\Signature{\Person{索斯特拉斯}，长老}\newline{}\Signature{\Person{阿摩尼乌斯}，长老}\newline{}\Signature{\Person{特翁}，长老}\newline{}\Signature{\Person{提兰努斯}，长老}\newline{}\Signature{\Person{波孔}，长老}\newline{}\Signature{\Person{科普雷斯}，长老}\newline{}\Signature{\Person{阿加图斯}，长老}\newline{}\Signature{\Person{阿摩纳斯}，长老}\newline{}\Signature{\Person{阿基利斯}，长老}\newline{}\Signature{\Person{奥瑞翁}，长老}\newline{}\Signature{\Person{保路斯}，长老}\newline{}\Signature{\Person{塞瑞努斯}，长老}\newline{}\Signature{\Person{塔莱拉乌斯}，长老}\newline{}\Signature{\Person{狄迪摩斯}，长老}\newline{}\Signature{\Person{狄奥尼修斯}，长老}\newline{}\Signature{\Person{赫拉克勒斯}，长老}\newline{}\Signature{\Person{塞拉皮翁}，执事}\newline{}\Signature{\Person{狄迪摩斯}，执事}\newline{}\Signature{\Person{普托拉瑞翁}，执事}\newline{}\Signature{\Person{犹斯图斯}，执事}\newline{}\Signature{\Person{塞拉斯}，执事}\newline{}\Signature{\Person{加约}，执事}\newline{}\Signature{\Person{狄迪摩斯}，执事}\newline{}\Signature{\Person{希厄拉克斯}，执事}\newline{}\Signature{\Person{德默特里乌斯}，执事}\newline{}\Signature{\Person{玛尔库斯}，执事}\newline{}\Signature{\Person{毛鲁斯}，执事}\newline{}\Signature{\Person{特奥纳斯}，执事}\newline{}\Signature{\Person{亚历山大}，执事}\newline{}\Signature{\Person{萨尔玛彤}，执事}\newline{}\Signature{\Person{玛尔库斯}，执事}\newline{}\Signature{\Person{卡尔蓬}，执事}\newline{}\Signature{\Person{科蒙}，执事}\newline{}\Signature{\Person{佐伊鲁斯}，执事}\newline{}\Signature{\Person{特里丰}，执事}\newline{}\Signature{\Person{阿摩尼乌斯}，执事}\par

% structure: p0032 source=h2 target=section reason=relative-heading-rank; base=h2

\section{三、信函}

\Person{亚历山大}，致\Place{亚历山大}城和\Place{马雷奥蒂斯}的司铎与执事，亲自向在场的你们，主内亲爱的弟兄们，谨致问候：\par

你们虽然已经迅速签署了我寄给那些跟随\Person{亚略}的人的信，敦促他们弃绝自己的不敬，并服从健全的\OriginalTerm{公教信仰}{Catholic faith}；并以此表明了你们的正统意向，以及与\OriginalTerm{天主教会}{Catholic Church}教义的一致；然而，因为我也曾就\Person{亚略}及其同伴的事，致信给各地的所有同工；我认为有必要召集你们这城市的圣职者，并也召唤你们\Place{马雷奥蒂斯}的圣职者；特别是因为你们当中的司铎\Person{卡雷斯}和\Person{皮斯图斯}，以及执事\Person{萨拉皮翁}、\Person{帕拉蒙}、\Person{佐西默}和\Person{伊勒内}，都已投奔\Person{亚略}一派，情愿与他们一同被罢黜；好使你们知道现在所写的事，并应宣布你们在这些事上的赞同，并为罢黜\Person{亚略}和\Person{皮斯图斯}一党投票。因为你们应当知道我所写的，并应各自如同亲自书写一般，存记在心。\par

% structure: p0035 source=h2 target=section reason=relative-heading-rank; base=h2

\section{四、致\Place{基诺波利斯}主教\Person{埃格隆}的信函——驳斥亚略派}

摘自圣\Person{亚历山大}，\Place{亚历山大}主教，致\Place{基诺波利斯}主教\Person{埃格隆}，驳斥亚略派的信函。\par

1. \OriginalTerm{自然意志}{natural will}是每一个理智本性所具有的自由官能，在其本质方面没有任何非自愿的成分。\par

2. \OriginalTerm{自然运作}{natural operation}是所有实体的内在运动。自然运作是每一种本性的实体性和\OriginalTerm{说明性的原理}{notifying reason}。自然运作是每一种实体的\OriginalTerm{表明特性的德能}{notifying virtue}。\par

% structure: p0039 source=h2 target=section reason=relative-heading-rank; base=h2

\section{五、论灵魂、身体与主的受难}

1. 圣言自天不吝赐下，是为浇灌我们的心田，只要我们的耳不是仅仅说话，而是聆听，便会为祂的德能预备好。因为正如雨水无地不能结实，同样，圣言无聆听不能结实，聆听无圣言也不可。再者，我们开口宣讲，圣言便结果实，同样，我们侧耳聆听，聆听也就结果实。所以，当圣言发挥其德能时，你们也要不吝地打开耳朵，前来聆听，并要除去一切恶意与不信。非常不良的两件事是恶意与不信，两者皆与\OriginalTerm{义德}{righteousness}相悖；因为恶意相反\OriginalTerm{爱德}{charity}，不信相反\OriginalTerm{信德}{faith}；正如苦与甜相反，黑暗与光明相反，恶与善相反，死亡与生命相反，虚谎与真理相反。因此，凡满溢这些与德性相斥的恶习之人，在某种方式上已是死人；因为恶毒者和不信者憎恨爱德与信德，而这样做的人就是\OriginalTerm{天主的仇敌}{enemies of God}。\par

2. 因此，亲爱的弟兄们，既然你们知道邪恶之徒和不信者是正义的仇敌，就当提防他们，持守信德和爱德；所有从世界之初直到今日的圣者，都是借着信德和爱德获得了救恩。你们要结出爱德的果实，不仅在言语上，也在行为上，即为天主的缘故，在一切敬虔的忍耐中彰显出来。看哪！主亲自向我们显示了祂的爱，不仅在言语上，也在行为上，因为祂舍弃了自己，作为我们救恩的代价。再者，我们受造，不像世界其余之物，只凭言语，也凭行为。因为天主以大能的一句话使世界存在，但祂创造我们，既靠祂言语的功效，也靠祂的作为。因为天主不单说：\BibleQuote{让我们照我们的肖像，按我们的模样造人} \BibleRef{创 1:26}，而且行为随言语而来；祂用地上的尘土造了人，使他符合自己的肖像和模样，并向人吹了一口生气，亚当就成了一个活的灵魂。\par

3. 但后来人因堕落而趋向死亡时，同一\OriginalTerm{造物主}{Artificer}就必须重新创造那形体，以获得救恩。因为那形体在地上腐烂；而那曾作为生命之气的生气，却与身体分离，被拘禁在黑暗之处，即所谓的\OriginalTerm{阴府}{Hades}。因此，灵魂与身体分开了；灵魂被放逐\Emph{\OriginalTerm{下到阴府}{ad inferos}}，而身体归于尘土；两者之间有了巨大的分离间隔；因为身体因血肉的解体而朽坏；灵魂既脱离了身体，其活动便停止了。正如国王被锁链捆绑，城邑就败落；或如将军被俘，军队便溃散；或如舵手被抛下，船只便沉没；同样，当灵魂被锁链捆锁，其身体就破碎；如同没有国王的城邑，身体的各部分就瓦解；如同失去将军的军队，他们就沉沦于死亡；正像失去舵手的船只那样。因此，只要身体还在，灵魂就治理人；正如国王治理城邑、将军治理军队、舵手治理船只。但从灵魂不可移动地与身体相连，并陷入错谬之时起，它便无力治理了；因此，它偏离了正途，追随诱惑者，沉溺于邪淫、偶像崇拜和流血；借着这些恶行，它毁灭了真正的人性。不单如此，它自己最终也被带到阴间，在那里被邪恶的诱惑者拘禁。否则，正如国王重建败落的城邑、将军重新召集溃散的军队、水手修复破损的船只，同样，我要说，在身体尚未归于尘土、自身也未被紧紧捆绑之前，灵魂曾向身体供应所需。但当灵魂被捆绑，不是被物质的锁链，而是被罪捆绑，因而无力行动时，它便将身体留在地上，自己被抛弃到阴间，成为死亡的脚凳，被众人藐视。\par

4. 人从乐园出来，进到一个充满不义、邪淫、奸淫和残忍凶杀的地方。在那里他找到了自己的毁灭；因为一切都串通起来使他死亡，使刚进入此地的人沦于丧亡。同时，人需要安慰、帮助和安息。因为人什么时候安好呢？在母胎中吗？但那时他被封闭里面，与死人几乎没有差别。当他由母乳喂养时吗？确实，连那时他也感觉不到任何喜乐。还是在他长大成人的时候呢？但那时，青春的情欲尤其给他带来危险。最后，是在他衰老的时候吗？不，那时他却开始呻吟，被衰老的重担和死亡的等待所压着。因为老年若不是等待死亡，又是什么呢？实在，地上所有居民都要死亡，青年人和老年人，孩童和成人，因为没有一个年龄或身形能免于死亡。那么，人为何被这极度的忧伤折磨呢？无疑，正是死亡的景象产生了哀伤；因为我们在死人身上看见面貌改变，形象死沉，身体因消瘦而萎缩，口哑，皮肤冰冷，尸身倒在地上，眼睛凹陷，四肢不动，肌肉消损，血脉凝结，骨头变白，关节脱离，全身化为尘土，人不再存在。那么，人是什么？我说，人是一朵短暂的花；在母胎中不显现，在青春时盛开，但在老年时枯萎，在死亡中逝去。\par

5. 然而，在经历了这一切对死亡的奴役和人性败坏之后，天主眷顾了祂按自己的肖像和模样所造的受造物；祂这样做，是为了不让它永远成为死亡的玩物。因此，天主从天上派遣祂\OriginalTerm{无躯体的圣子}{incorporeal Son}，在童贞女的胎中取肉躯；这样，祂就完全像你们一样，成了人；为拯救迷失的人，并聚集祂所有分散的肢体。因为基督将人性结合于自己的\OriginalTerm{位格}{person}时，就把死亡因身体分离而分散的重新聚合起来。基督受难，是为叫我们永远生活。\par

否则基督为何应当死亡？祂曾犯有何等死罪？祂本已披戴光荣，为何穿上肉躯？祂既是天主，为何成为人？祂既在天为王，为何降来尘世，在童贞的胎中降生成人？我要问，有何必要迫使天主降来尘世，取肉躯，被裹以襁褓，卧于马槽摇篮，被哺以母乳，由仆人手中接受圣洗，被高举在十字架上，被葬于尘世坟墓，第三日从死者中复活？我说，有何必要迫使祂如此？显然祂是为了人而忍受羞辱，为解脱人的死亡；祂又曾像先知的话那样呼喊：\BibleQuote{我忍受，犹如生产的妇女。}\BibleRef{依 42:14} 确实，祂为了我们承受了忧苦、羞辱、磨难，甚至死亡本身和埋葬。因为祂自己藉先知这样说：\BibleQuote{我沉入深渊。}\BibleRef{纳 2:4} 谁使祂如此沉落？是不虔敬的民众。看，世人，看\Person{以色列}怎样回报了祂！她杀害了她的恩主，以恶报善，以忧苦报喜悦，以死亡报生命。他们将祂钉在\OriginalTerm{树}{tree}上杀害了，而祂曾复活他们的死者，治愈他们的残者，洁净他们的癞病人，赐光明给他们的瞎子。看，世人！看，众百姓，这些新的奇事！他们将那铺展大地的，悬于\OriginalTerm{树}{tree}上；他们将那奠定世界基础的，用钉穿透；他们将那环绕诸天的，加以限制；他们将那赦免罪人的，予以捆绑；他们给那使他们饮于义德的，喝醋；他们使那赐予他们生命之粮的，吃胆；他们使那治愈他们手足的，手和脚遭受朽坏；他们强行闭合了那赐予他们恢复视力的，眼睛；他们将那在祂受难前及在祂悬于\OriginalTerm{树}{tree}上时复活他们死者生命的，交付于坟墓。\par

6. 因为当我们的主在十字架上受苦时，坟墓裂开，地狱显露，灵魂跃出，死者复活，其中许多人显现于\Place{耶路撒冷}，正当十字架的奥迹完成之时；那时我们的主踏碎了死亡，化解了仇恨，捆绑了壮士，举起了十字架的胜利标记，祂的身体被高举其上，为使身体显于高处，而死亡被践踏在肉躯的脚下。那时天上的异能惊奇，天使惊讶，元素颤栗，所有受造物震动，因他们目睹这新奥迹，以及宇宙中正在上演的可怖景象。然而全体民众，仿佛不认识这奥迹，嘲弄基督而欢跃；尽管大地摇撼，山岳、山谷和海洋震动，天主的所有受造物都惊惶失措。天上的光体惊惧，太阳逃遁，月亮隐没，星辰收敛光芒，白昼终结；天使在惊讶中离开\Place{圣殿}，在帐幔裂开之后，黑暗覆盖了大地，其主闭上了双眼。同时，\Place{地狱}充满光辉，因那异星降了下来。主确实不以肉身，而以神魂降入\Place{地狱}。祂其实在处处工作，因为当祂以肉身复活死者时，祂以神魂解放了他们的灵魂。因为当主的身体悬在十字架上时，正如我们所说，坟墓开启了；\Place{地狱}的门闩被除去，死者得回生命，灵魂被遣返世界，这一切是因为主战胜了\Place{地狱}，踏碎了死亡，使仇敌蒙羞；因此，灵魂从\Place{地狱}出来，死者显现于地上。\par

7. 所以你们看，基督之死的效能何其伟大，因为没有任何受造物以平静的心承受祂的倒下，诸元素也未平静承受祂的苦难，大地不能留住祂的肉身，\Place{地狱}不能扣留祂的神魂。在基督的苦难中，万物被扰乱、震撼。主呼唤，如同从前对\Person{拉匝禄}那样：\Quote{出来吧，死者，从你们的坟墓和隐密处出来；因为我基督赐给你们复活。} 因为那时，大地不能长久保持安葬于其中的主之肉身；她呼喊：\Quote{哦，我主，赦免我的罪过，救我免于你的义怒，解除我的诅咒，因为我接受了义人的血，却未曾掩盖世人的身体，也未掩盖你自己的身体！} 这奇妙奥迹究竟为何？主啊，你为什么降来尘世，若不是为了那散布各处的人？因为在各处，你美好的肖像已被散播。不！但只要你发出一言，瞬间所有身体都会站在你面前。如今，你既已来到世上，寻找你所塑造的肢体，就为你自己的人负责吧，接受那托付于你的，收回你的肖像，你的\Person{亚当}。然后，主在死后第三日复活了，从而引领人认识天主圣三。于是人类各邦都由基督得救了。一人承受了审判，无数人获得了赦免。此外，祂成为相似祂所拯救的人，升到了高天，在祂的圣父面前奉献，不是金银宝石，而是祂按照自己的肖像和模样所造的人；圣父举起祂在祂的右边，使祂坐在高高的宝座上，立祂为万民的审判者，天使军旅的统帅，革鲁宾的驭者，真\Place{耶路撒冷}的子嗣，童贞的净配，并永远为君王。阿们。\par

% structure: p0048 source=h2 target=section reason=relative-heading-rank; base=h2

\section{六、抄本中的增补，附异文}

天主因此愿意眷顾祂依照自己的肖像和模样所塑造的\OriginalTerm{形式}{form}，在这些末后的日子里，派遣了祂的\OriginalTerm{无形体的}{incorporeal}独生子进入世界；祂在童贞女胎中\OriginalTerm{降生成人}{incarnate}，诞生为完全的人，为将堕落的人扶起，重新聚集祂四散的肢体。因为基督为什么应该死去呢？祂是被控犯有死罪吗？既然祂是天主，为什么成了人？那在诸天之上为王的，为什么降临地上？是谁催迫天主降临地上，由圣洁童贞女取血肉，裹以襁褓，卧于马槽，哺乳喂养，在\Place{约但河}受洗，受民众戏弄，被钉于木架，葬于地腹，第三日从死者中复活；为了救赎，以命换命，以血换血，以死代死？因为基督通过死亡，清偿了人因死亡所负的罪债。啊，这新的、\OriginalTerm{不可言喻的}{ineffable}奥迹！审判者受了审判；赦免罪过的那位被捆绑；曾创造世界的那位受了戏弄；伸展诸天的那位被悬于十字架上；以玛纳为食粮的那位被喂以苦胆；赐生命的那位死了；复活死者的那位被交付于坟墓。诸掌权者惊愕，天使惊奇，元素战栗，整个受造宇宙震动，大地颤栗，地基摇动；太阳遁去，元素颠覆，白昼之光消退；因为它们不忍注视被钉的主。受造物惊异地说：\Quote{这是何等新奇的奥迹？审判者受审而缄默；不可见者被见而不羞愧；不可把握者被捉住而不恼怒；不可度量者被容于度量中而不抗拒；\OriginalTerm{无感受的}{impassable}受难而不报复己伤；不死者死去而无怨言；属天者被埋葬而心平气和。} 我说：这是何等奥迹？受造物确实惊愕得目瞪口呆。但当我们的主从死亡中复活并践踏了死亡，当祂束缚了壮士而释放了人，那时一切受造物都惊奇：为了\Person{亚当}的缘故受审的审判者，被看见的不可见者，受难的无感受者，死去的不死者，被埋在地里的属天者。因为我们的主成了人；祂被定罪为施怜悯；祂被捆绑为释放人；祂被捉拿为解救人；祂受苦为医治我们的痛苦；祂死去为恢复我们的生命；祂被埋葬为使我们复活起来。因为当我们的主受苦时，是祂的人性受苦，就是祂与人相似的（部分）；祂化解了与祂相似之人的痛苦，并通过死亡摧毁了死亡。正是为此，祂降临地上，为要藉着追逐死亡，杀死那杀人的叛逆者。因为一人受了审判，千万人得释放；一人被埋葬，千万人复活。祂是天主与人之间的\OriginalTerm{中保}{Mediator}；祂是所有人的复活与救恩；祂是迷途者的向导，已获自由者的牧者，死者的生命，\OriginalTerm{革鲁宾}{cherubim}的战车御者，天使的旌旗持有者，万王之王，愿光荣归于祂，直到世世代代。阿们。\par

\SourceRecord{Translated by James B.H. Hawkins.；Ante-Nicene Fathers；Vol. 6.；Edited by Alexander Roberts, James Donaldson, and A. Cleveland Coxe.；Buffalo, NY: Christian Literature Publishing Co.,；1886.；<http://www.newadvent.org/fathers/0622.htm>.}
